\documentclass[11pt]{article}

\usepackage[utf8]{inputenc}
\usepackage[T1]{fontenc}
\usepackage[margin=1in]{geometry}
\usepackage{amsmath,amssymb}
\usepackage{booktabs}
\usepackage{enumitem}
\usepackage[hidelinks]{hyperref}

\newcommand{\R}{\mathbb{R}}
\newcommand{\N}{\mathbb{N}}
\newcommand{\catProg}{\mathbf{Prog}}
\newcommand{\catAuth}{\mathbf{Auth}}

\title{\bfseries MANIFOLD\\[2mm]\Large A Mathematical Engine for Mapping Vulnerabilities}
\author{MANIFOLD Project}
\date{Whitepaper v0.1 --- working draft, not peer reviewed}

\begin{document}
\maketitle

\begin{abstract}
\noindent MANIFOLD treats a software artifact not as a bag of rules to match, but
as a \emph{mathematical object}: a typed, weighted graph enriched with
\emph{algebraic}, \emph{spectral}, \emph{topological} and \emph{geometric}
structure. Over this object we compute a scalar field of \textbf{vulnerability
potential} $V(x)$, which surfaces anomalous regions as geometric / topological /
spectral \emph{features} rather than as literal string matches. An
\textbf{autonomous LLM agent} navigates the resulting ``manifold'', formulates
hypotheses about those regions, and dispatches \textbf{formal verifiers} (SMT,
symbolic execution, abstract interpretation) that confirm or refute each
hypothesis --- closing the loop between statistical intuition and mathematical
proof. This document defines the intermediate representation, formalizes each
mathematical layer, states the mapping laws that connect mathematical features
to vulnerability classes, and specifies the agent architecture and roadmap.
\end{abstract}

\section{Motivation}
Traditional scanners (SAST/DAST) enumerate syntactic patterns. They are brittle:
they miss novel vulnerabilities and drown analysts in false positives. Two
capabilities are missing:
\begin{enumerate}[label=(\arabic*)]
  \item A \emph{unified, quantitative representation} of ``program structure''
        onto which many \emph{independent} signals can be projected and fused.
  \item A \emph{reasoning loop} that converts ``suspicious region'' into
        ``proven vulnerability'' --- the LLM supplies the intuition, mathematics
        supplies the proof.
\end{enumerate}
MANIFOLD provides both.

\section{The core idea}
\[
\text{artifact}\ \xrightarrow{\quad}\ \text{IR (typed graph)}
\ \xrightarrow{\quad}\ \text{mathematical layers}
\ \xrightarrow{\quad}\ \text{scalar field } V(x)
\]
A vulnerability is \emph{not} a pattern; it is a \emph{violation of a structural
invariant} (a broken naturality condition, a persistent cycle, a negative
curvature chokepoint, a taint flow across a spectral trust cut). Each layer
detects a different \emph{kind} of invariant violation, and the fused field
$V(x)$ ranks regions for the agent to inspect.

\section{The Intermediate Representation (IR)}
\begin{quote}
\textbf{Definition 1 (Program graph).} A \emph{program graph} is a tuple
$G = (V, E, \tau_V, \tau_E, \omega)$ where
\begin{itemize}[leftmargin=1.5em]
  \item $V$ is a finite set of \emph{nodes} (modules, functions, blocks,
        statements, calls, variables, sources, sinks);
  \item $E \subseteq V \times V$ is a set of directed \emph{edges};
  \item $\tau_V : V \to \N_K$ assigns each node a \emph{kind} from
        $K = \{\text{module, function, block, statement, call, assign, variable,
        parameter, source, sink}\}$;
  \item $\tau_E : E \to \mathcal{T}$ assigns each edge a \emph{type}
        $\mathcal{T} = \{\text{control, data, call, taint, trust, auth}\}$;
  \item $\omega : E \to \R_{\ge 0}$ assigns each edge a non-negative
        \emph{weight} (default $1$).
\end{itemize}
\end{quote}
The IR is the single source of truth. Every adapter (source code, binary, web
API, LLM agent) must produce a program graph; every analysis consumes it. The
schema is language-agnostic and JSON-serializable (the Python and Rust layers
share the identical schema).

\section{Mathematical layers}
For each layer we fix a subgraph by edge type $\tau_E$, then compute structure.

\subsection{Spectral layer (algebraic graph theory)}
For an edge type $t$, define the symmetric \emph{adjacency} $A_t$ and
\emph{degree} $D_t$, and the \emph{combinatorial Laplacian}
\[
L_t = D_t - A_t .
\]
$L_t$ is positive semidefinite; its spectrum
$0 = \lambda_1 \le \lambda_2 \le \cdots$ encodes global structure:
\begin{itemize}[leftmargin=1.5em]
  \item $\lambda_2$ is the \emph{algebraic connectivity} (Fiedler value). Its
        eigenvector $f_2$ --- the \emph{Fiedler vector} --- induces a canonical
        cut $S^+ = \{v : f_2(v) \ge 0\}$.
  \item The \emph{spectral embedding}
        $\Phi^{(d)}(v) = (f_2(v), f_3(v), \dots, f_{d+1}(v))$ maps nodes into
        $\R^d$.
  \item The Cheeger inequality bounds the bottleneck ratio $h(G)$ by
        $\tfrac{\lambda_2}{2} \le h(G) \le \sqrt{2\,d_{\max}\lambda_2}$.
\end{itemize}
\textbf{Signal 1 (spectral anomaly).} The reconstruction error of a node under a
low-rank eigenmap reconstruction, and its distance from the Fiedler cut, are
used as anomaly scores.

\subsection{Topological layer (persistent homology)}
From the program graph we build a \emph{filtration} of simplicial complexes
(e.g.\ a Vietoris--Rips or flag complex over the shortest-path metric, or a
filtration driven by $\omega$). Persistent homology yields, for each dimension
$k$, a \emph{persistence diagram} $\mathrm{Dgm}_k$ of birth--death pairs.
$\beta_0$ counts components, $\beta_1$ counts independent cycles; points far
from the diagonal have high persistence. The \textbf{Mapper} algorithm
(Singh--M\'emoli--Carlsson) builds a navigable \emph{nerve} from a filter
function $f: V \to \R$, producing the map that gives the project its name.

\subsection{Geometric layer (discrete curvature)}
The \emph{Ollivier--Ricci curvature} of an edge $(u,v)$ is
\[
\kappa(u,v) = 1 - \frac{W_1(m_u, m_v)}{d(u,v)},
\]
with lazy random-walk measures $m_u, m_v$ and Wasserstein-1 distance $W_1$.
Negative curvature marks \emph{funnels} through which many shortest paths pass.
The \emph{Forman--Ricci} curvature (unweighted) is $\kappa_F(u,v) = 4 - \deg(u)
- \deg(v)$, cheaply flagging high-degree hubs.

\subsection{Algebraic layer (lattices \& category theory)}
\emph{Abstract interpretation.} The \emph{taint lattice}
$\mathbb{T} = \{\bot = \text{clean} \sqsubseteq \top = \text{tainted}\}$ (a
product lattice over tags) gives a \emph{Galois connection}
$(\alpha, \gamma)$ between the concrete collecting semantics and the abstract
taint domain, guaranteeing soundness. \emph{Category theory.} The program is a
category $\catProg$; effects are monads, authorization is a functor
$F : \catProg \to \catAuth$. \textbf{Design law (naturality of taint):} a
vulnerability is a taint flow that breaks the naturality square of $F$.

\subsection{Formal layer (symbolic execution \& model checking)}
Symbolic execution computes a path condition $\Phi$ and symbolic state $\sigma$;
an SMT solver checks $\mathrm{SAT}(\Phi \wedge \text{bad})$, whose model is a
concrete exploit. Reachability is the CTL formula $\mathbf{EF}\,\mathrm{bad}$;
taint propagation is a finite automaton whose accepted language is exactly the
set of source $\to$ sink paths.

\section{The vulnerability scalar field}
\[
V(x) = \sigma\!\Big(\textstyle\sum_i \alpha_i \, s_i(x)\Big),
\]
where $s_i$ are normalized layer signals (spectral anomaly, persistence
outlier, curvature, lattice height, formal reachability), $\alpha_i$ are
(learnable) weights, and $\sigma$ a logistic squeeze. $V$ is a ranking, a
Mapper filter, and a cost field for navigation.

\section{Mapping laws}
\begin{center}
\begin{tabular}{cll}
\toprule
\# & Law & Vulnerability class \\
\midrule
L1 & Chokepoint ($\kappa \ll 0$) & privilege escalation, auth bypass \\
L2 & Cycle (persistent $H_1$) & reentrancy, infinite recursion \\
L3 & Trust cut (Fiedler cut) & injection across trust boundary \\
L4 & Naturality break & arbitrary taint violation \\
L5 & Persistence outlier & real vs.\ spurious features \\
L6 & Reachability ($\mathrm{SAT}(\Phi \wedge \text{bad})$) & concrete exploit \\
\bottomrule
\end{tabular}
\end{center}

\section{The autonomous LLM agent}
A closed loop: (1) \emph{ingest} into the IR; (2) \emph{embed \& map} ($V(x)$,
spectral embedding, persistence, curvature, Mapper); (3) \emph{hypothesize}
(the LLM proposes candidates from the manifold); (4) \emph{verify} (the formal
layer proves or refutes); (5) \emph{refine} (findings re-weight and re-embed);
(6) \emph{report}. The LLM provider layer targets any OpenAI-compatible endpoint
via \textbf{LiteLLM}. The LLM does not do raw detection: mathematics proposes
regions, the LLM proposes hypotheses, the formal layer proposes proofs.

\section{Implementation status}
\begin{center}
\begin{tabular}{lll}
\toprule
Phase & Deliverable & Status \\
\midrule
F0 & Whitepaper, architecture & done \\
F1 & IR + Python SAST ingest + taint & done \\
F2 & Spectral + geometric kernels (Rust) & done \\
F3 & TDA (persistent homology, Mapper) & done \\
F4 & Formal verification (Z3, angr) & pending \\
F5 & Autonomous LLM agent (LiteLLM) & pending \\
F6 & Visualization + end-to-end demo & pending \\
F7 & Binary / web / LLM adapters & pending \\
\bottomrule
\end{tabular}
\end{center}

\section{Validation and responsible use}
Validation uses toy micro-benchmarks with known ground truth, then labeled
corpora (OWASP Benchmark, SARD) once F4 lands, measuring precision/recall of the
fused field versus raw taint. MANIFOLD is a defensive / authorized-testing tool;
we release under the assumption of authorized use and encourage coordinated
disclosure.

\begin{thebibliography}{9}
\item M.~Fiedler, \emph{Algebraic connectivity of graphs}, 1973.
\item F.~Chung, \emph{Spectral Graph Theory}, 1997.
\item Edelsbrunner, Letscher, Zomorodian, \emph{Topological persistence and simplification}, 2002.
\item Singh, M\'emoli, Carlsson, \emph{Topological Methods for the Analysis of High Dimensional Data Sets}, 2007.
\item Y.~Ollivier, \emph{Ricci curvature of Markov chains on metric spaces}, 2009.
\item P.~Cousot \& R.~Cousot, \emph{Abstract interpretation: a unified lattice model}, 1977.
\item S.~Mac~Lane, \emph{Categories for the Working Mathematician}, 1971.
\item J.~King, \emph{Symbolic execution and program testing}, 1976.
\item E.~M.~Clarke et al., \emph{Automatic verification of finite-state concurrent systems}, 1986.
\end{thebibliography}

\end{document}
